% Options for packages loaded elsewhere
\PassOptionsToPackage{unicode}{hyperref}
\PassOptionsToPackage{hyphens}{url}
\PassOptionsToPackage{dvipsnames,svgnames,x11names}{xcolor}
\documentclass[
  10pt,
  fontset=fandol]{ctexart}
\usepackage{xcolor}
\usepackage[margin=16mm]{geometry}
\usepackage{amsmath,amssymb}
\setcounter{secnumdepth}{-\maxdimen} % remove section numbering
\usepackage{iftex}
\ifPDFTeX
  \usepackage[T1]{fontenc}
  \usepackage[utf8]{inputenc}
  \usepackage{textcomp} % provide euro and other symbols
\else % if luatex or xetex
  \usepackage{unicode-math} % this also loads fontspec
  \defaultfontfeatures{Scale=MatchLowercase}
  \defaultfontfeatures[\rmfamily]{Ligatures=TeX,Scale=1}
\fi
\usepackage{lmodern}
\ifPDFTeX\else
  % xetex/luatex font selection
\fi
% Use upquote if available, for straight quotes in verbatim environments
\IfFileExists{upquote.sty}{\usepackage{upquote}}{}
\IfFileExists{microtype.sty}{% use microtype if available
  \usepackage[]{microtype}
  \UseMicrotypeSet[protrusion]{basicmath} % disable protrusion for tt fonts
}{}
\makeatletter
\@ifundefined{KOMAClassName}{% if non-KOMA class
  \IfFileExists{parskip.sty}{%
    \usepackage{parskip}
  }{% else
    \setlength{\parindent}{0pt}
    \setlength{\parskip}{6pt plus 2pt minus 1pt}}
}{% if KOMA class
  \KOMAoptions{parskip=half}}
\makeatother
\usepackage{longtable,booktabs,array}
\newcounter{none} % for unnumbered tables
\usepackage{calc} % for calculating minipage widths
% Correct order of tables after \paragraph or \subparagraph
\usepackage{etoolbox}
\makeatletter
\patchcmd\longtable{\par}{\if@noskipsec\mbox{}\fi\par}{}{}
\makeatother
% Allow footnotes in longtable head/foot
\IfFileExists{footnotehyper.sty}{\usepackage{footnotehyper}}{\usepackage{footnote}}
\makesavenoteenv{longtable}
\setlength{\emergencystretch}{3em} % prevent overfull lines
\providecommand{\tightlist}{%
  \setlength{\itemsep}{0pt}\setlength{\parskip}{0pt}}
\usepackage{amsmath,amssymb,mathtools,needspace}
\xeCJKDeclareCharClass{CJK}{"2032,"0400->"04FF,"2160->"216B,"2460->"2473,"25A0,"25C7,"25CB,"25CF}
\IfFontExistsTF{Arial Unicode MS}{\xeCJKsetup{AutoFallBack=true}\setCJKfallbackfamilyfont{\CJKrmdefault}{Arial Unicode MS}}{}
\setcounter{secnumdepth}{0}
\setlength{\emergencystretch}{2em}
\allowdisplaybreaks
\usepackage{bookmark}
\IfFileExists{xurl.sty}{\usepackage{xurl}}{} % add URL line breaks if available
\urlstyle{same}
\hypersetup{
  pdftitle={误差来源与误差分析的重要性},
  colorlinks=true,
  linkcolor={Maroon},
  filecolor={Maroon},
  citecolor={Blue},
  urlcolor={Blue},
  pdfcreator={LaTeX via pandoc}}

\title{误差来源与误差分析的重要性}
\author{}
\date{}

\begin{document}
\maketitle

\subsubsection{1.2 误差来源与误差分析的重要性}\label{ux8befux5deeux6765ux6e90ux4e0eux8befux5deeux5206ux6790ux7684ux91cdux8981ux6027}

用计算机解决科学计算问题首先要建立数学模型,它是对被描述的实际问题进行抽象、简化而得到的,因而是近似的.我们把数学模型与实际问题之间出现的这种误差称为模型误差.只有实际问题提法正确,建立数学模型时又抽象、简化得合理,才能得到好的结果.由于这种误差难以用数量表示,通常都假定数学模型是合理的,这种误差可忽略不计,在``数值分析''中不予讨论.在数学模型中往往还有一些根据观测得到的物理量,如温度、长度、电压等,这些参量显然也包含误差.这种由观测产生的误差称为观测误差,在``数值分析''中也不讨论这种误差.数值分析只研究用数值方法求解数学模型产生的误差.

当数学模型不能得到精确解时,通常要用数值方法求它的近似解,其近似解与精确解之间的误差称为截断误差或方法误差.例如,当函数 \(f(x)\) 用 Taylor多项式

\[
P_n(x)=f(0)+\frac{f'(0)}{1!}x+\frac{f''(0)}{2!}x^2+\cdots+\frac{f^{(n)}(0)}{n!}x^n
\]

近似代替时,数值方法的截断误差是

\[
R_n(x)=f(x)-P_n(x)=\frac{f^{(n+1)}(\xi)}{(n+1)!}x^{n+1},\qquad \xi\text{ 在 }x\text{ 与 }0\text{ 之间}.
\]

\href{../../source-images/pdf-016.jpeg}{查看原页}

有了求解数学问题的计算公式以后，用计算机进行数值计算时，由于计算机的字长有限，原始数据在计算机上表示会产生误差，计算过程又可能产生新的误差，这种误差称为舍入误差，例如，用3.14159近似代替\(\pi\)，产生的误差

\[
R = \pi - 3.14159 = 0.0000026\cdots
\]

就是舍入误差. 在``数值分析''中除了研究数学问题的算法外，还要研究计算结果的误差是否满足精度要求，这就是误差估计问题. 本书主要讨论算法的截断误差与舍入误差，对舍入误差通常只作一些定性分析. 下面举例说明误差分析的重要性.

\textbf{例1.1} 计算 \(I_n = \mathrm{e}^{-1} \int_{0}^{1} x^n \mathrm{e}^{x} \mathrm{d}x (n = 0,1,\cdots)\)，并估计误差.

\textbf{解} 由分部积分可得计算 \(I_n\) 的递推公式

\[
I_n = 1 - n I_{n-1} \quad (n = 1,2,\cdots), \tag{1.2.1}
\]

\[
I_0 = \mathrm{e}^{-1} \int_{0}^{1} \mathrm{e}^{x} \mathrm{d}x = 1 - \mathrm{e}^{-1}.
\]

若计算出 \(I_0\)，代入式(1.2.1)，可逐次求出 \(I_1, I_2, \cdots\) 的值. 要算出 \(I_0\) 就要先计算 \(\mathrm{e}^{-1}\)，若用 Taylor多项式展开部分和

\[
\mathrm{e}^{-1} \approx 1 + (-1) + \frac{(-1)^2}{2!} + \cdots + \frac{(-1)^k}{k!},
\]

并取 \(k=7\)，用四位小数计算，则得 \(\mathrm{e}^{-1} \approx 0.3679\)，截断误差

\[
R_7 = |\mathrm{e}^{-1} - 0.3679| \leqslant \frac{1}{8!} < \frac{1}{4} \times 10^{-4}.
\]

计算过程中小数点后第五位的数字按四舍五入原则舍入，由此产生的舍入误差这里先不讨论. 当初值取为 \(I_0 \approx 0.6321 = \tilde{I}_0\) 时，用式(1.2.1)递推的计算公式为 方案(A)

\[
\begin{cases}
\tilde{I}_0 = 0.6321, \\
\tilde{I}_n = 1 - n \tilde{I}_{n-1} \quad (n = 1,2,\cdots).
\end{cases}
\]

计算结果如表1.1的 \(\tilde{I}_n\) 列所示. 用 \(\tilde{I}_0\) 近似 \(I_0\) 产生的误差 \(E_0 = I_0 - \tilde{I}_0\) 就是初值误差，它对后面计算结果是有影响的.

表1.1

{\def\LTcaptype{none} % do not increment counter
\begin{longtable}[]{@{}
  >{\raggedright\arraybackslash}p{(\linewidth - 10\tabcolsep) * \real{0.1667}}
  >{\raggedright\arraybackslash}p{(\linewidth - 10\tabcolsep) * \real{0.1667}}
  >{\raggedright\arraybackslash}p{(\linewidth - 10\tabcolsep) * \real{0.1667}}
  >{\raggedright\arraybackslash}p{(\linewidth - 10\tabcolsep) * \real{0.1667}}
  >{\raggedright\arraybackslash}p{(\linewidth - 10\tabcolsep) * \real{0.1667}}
  >{\raggedright\arraybackslash}p{(\linewidth - 10\tabcolsep) * \real{0.1667}}@{}}
\toprule\noalign{}
\begin{minipage}[b]{\linewidth}\raggedright
\(n\)
\end{minipage} & \begin{minipage}[b]{\linewidth}\raggedright
\(\tilde{I}_n\) (用方案(A)计算)
\end{minipage} & \begin{minipage}[b]{\linewidth}\raggedright
\(I_n^*\) (用方案(B)计算)
\end{minipage} & \begin{minipage}[b]{\linewidth}\raggedright
\(n\)
\end{minipage} & \begin{minipage}[b]{\linewidth}\raggedright
\(\tilde{I}_n\) (用方案(A)计算)
\end{minipage} & \begin{minipage}[b]{\linewidth}\raggedright
\(I_n^*\) (用方案(B)计算)
\end{minipage} \\
\midrule\noalign{}
\endhead
\bottomrule\noalign{}
\endlastfoot
0 & 0.6321 & 0.6321 & 5 & 0.1480 & 0.1455 \\
1 & 0.3679 & 0.3679 & 6 & 0.1120 & 0.1268 \\
2 & 0.2642 & 0.2643 & 7 & 0.2160 & 0.1121 \\
3 & 0.2074 & 0.2073 & 8 & -0.728 & 0.1035 \\
4 & 0.1704 & 0.1708 & 9 & 7.552 & 0.0684 \\
\end{longtable}
}

\href{../../source-images/pdf-017.jpeg}{查看原页}

从表1.1可以看到，\(\tilde{I}_8\) 出现负值，这与一切 \(I_n > 0\) 相矛盾. 实际上，由积分估值得

\[
\frac{e^{-1}}{n+1} = e^{-1}\left(\min_{0 \leqslant x \leqslant 1} e^x\right)\int_0^1 x^n \mathrm{d}x < I_n < e^{-1}\left(\max_{0 \leqslant x \leqslant 1} e^x\right)\int_0^1 x^n \mathrm{d}x = \frac{1}{n+1}. \tag{1.2.2}
\]

因此，当 \(n\) 较大时，用 \(\tilde{I}_n\) 近似 \(I_n\) 显然是不正确的. 这里，计算公式与每步计算都是正确的，那么，是什么原因使计算结果错误呢？主要就是初值 \(\tilde{I}_0\) 有误差 \(E_0 = I_0 - \tilde{I}_0\)， 由此引起以后各步计算的误差 \(E_n = I_n - \tilde{I}_n\) 满足关系 \(E_n = -nE_{n-1}(n=1,2,\cdots)\). 容易推得

\[
E_n = (-1)^n n! E_0,
\]

这说明 \(\tilde{I}_0\) 有误差 \(E_0\)，则 \(\tilde{I}_n\) 就是 \(E_0\) 的 \(n!\) 倍误差. 例如，\(n=8\)，若 \(|E_0| = \frac{1}{2} \times 10^{-4}\)，则 \(|E_8| = 8! \times |E_0| > 2\). 这就说明 \(\tilde{I}_8\) 完全不能近似 \(I_8\) 了. 我们现在换一种计算方案. 由式(1.2.2)取 \(n=9\)，得

\[
\frac{e^{-1}}{10} < I_9 < \frac{1}{10},
\]

粗略取 \(I_9 \approx \frac{1}{2}\left(\frac{1}{10} + \frac{e^{-1}}{10}\right) = 0.0684 = I_9^*\)， 然后将式(1.2.1)倒过来算，即由 \(I_9^*\) 算出 \(I_8^*, I_7^*, \cdots, I_1^*\)，公式为 方案(B)

\[
\begin{cases} I_9^* = 0.0684, \\ I_{n-1}^* = \frac{1}{n}(1 - I_n^*) \quad (n = 9, 8, \cdots, 1). \end{cases}
\]

计算结果如表1.1中的 \(I_n^*\) 列所示. 可以发现，\(I_0^*\) 与 \(I_0\) 的误差不超过 \(10^{-4}\). 由于 \(|E_0^*| = \frac{1}{n!}|E_n^*|, E_0^*\) 比 \(E_n^*\) 缩小了 \(n!\) 倍，因此，尽管 \(E_9^*\) 较大，但由于误差逐步缩小，故可用 \(I_n^*\) 近似 \(I_n\). 反之，当用方案(A)计算时，尽管初值 \(\tilde{I}_0\) 相当准确，但由于误差传播是逐步扩大的，因而计算结果不可靠. 此例说明，在数值计算中如不注意误差分析，用了类似于方案(A)的计算公式，就会出现``差之毫厘，失之千里''的错误结果. 尽管数值计算中估计误差比较困难，但仍应重视计算过程中的误差分析.

\begin{center}\rule{0.5\linewidth}{0.5pt}\end{center}

\subsection{相关原页校注（agent 补充）}\label{ux76f8ux5173ux539fux9875ux6821ux6ce8agent-ux8865ux5145}

以下校注来自本节覆盖的原页，可能也涉及同页相邻小节；原文照录与校正说明分开保留。

\subsubsection{原书 PDF 页 17 的校注}\label{ux539fux4e66-pdf-ux9875-17-ux7684ux6821ux6ce8}

\href{../pages/pdf-017.md}{完整原页转录} · \href{../../source-images/pdf-017.jpeg}{原页图像}

校注记录：CH01-E001。

\begin{quote}
校注（agent 补充）：本页定义1.1末尾``绝''字与下一页``对误差''相接。
\end{quote}

\begin{quote}
校注（agent 补充，CH01-E001）：原书写``则 \(\tilde{I}_n\) 就是 \(E_0\) 的 \(n!\) 倍误差''，此处照录。由上式 \(E_n=(-1)^n n!E_0\) 可知，应理解为近似值 \(\tilde{I}_n\) 的误差绝对值为初值误差绝对值的 \(n!\) 倍；疑为原书表述不严谨，不是 OCR 错误。
\end{quote}

\subsection{本节覆盖原页的校注记录}\label{ux672cux8282ux8986ux76d6ux539fux9875ux7684ux6821ux6ce8ux8bb0ux5f55}

下列记录按原页关联，含同页相邻小节的校注；计算时同时读取上文校注和完整原页。

\begin{itemize}
\tightlist
\item
  \texttt{CH01-E001}：原书 PDF 页 17
\end{itemize}

\end{document}
